\section{Literature Review}
\label{sec:lit}

This review was revised after the Phase 1 evaluation, which asked for four
things the first version lacked: a distinction between survey-based perceived
quality and platform ratings, applied rather than foundational citations for the
model choices with a justification within the tree family, literature support
for the features, and an evaluation aligned with a team's limited capacity.
Each is addressed in the subsection it belongs to.

\subsection{Fulfilment, delivery promises and satisfaction}

The e-commerce service-quality literature is direct about which lever matters.
\citet{parasuraman2005esqual} decompose electronic service quality into
efficiency, \emph{fulfilment}, system availability and privacy, and find
fulfilment --- the accuracy of delivery promises and the condition in which
goods arrive --- among the strongest contributors to perceived quality;
\citet{rao2011elsq} sharpen this into logistics terms across 260 online
retailers. Both measure a survey construct. What a marketplace observes is a
rating, and ratings are not surveys: they are written by a self-selected
minority, J-shaped by acquisition and under-reporting biases
\citep{hu2009jshaped, hu2017selfselection}, and moved by product, price and
seller behaviour as much as by service \citep{dellarocas2003digitization}.
Transactional evidence bridges the two: across more than 15 million marketplace
orders, \citet{deshpande2023logistics} trace a chain --- logistics performance
drives the rating a customer leaves, and that rating moves seller sales. That
chain makes a low-review model a commercial instrument rather than a
diagnostic, and the rating is the construct this project predicts.

The promise itself is a second decision variable. \citet{cui2024promise} show
that committing to a faster delivery date raises sales but increases returns and
depresses retention, so an estimated delivery date is a business decision, not a
forecast, and should not be expected to behave like one. Prediction pays only
when it feeds a decision: \citet{kandula2021prescriptive} route machine-learned
delivery outcomes into a scheduling optimiser and report cost savings of up to
10.2\%, the pattern Table~\ref{tab:levers} follows. The same literature warns
what is \emph{not} observable in an order table --- product expectations, seller
communication, returns handling --- and any ceiling reported here must be read
against that omission.

\subsection{Model families for tabular transactional data}

Regularised linear and logistic models are the natural baseline for structured
data of this size: cheap, interpretable and, as Section~\ref{sec:results}
shows, revealing about distribution shift. Within the tree family the
literature knows three ways to improve on one tree: a random forest bags
decorrelated trees to cut variance \citep{breiman2001rf}; boosting fits trees to
residuals with shrinkage to cut bias \citep{friedman2001gbm}. Across the large
empirical comparisons --- ten algorithms on eleven problems
\citep{caruana2006empirical}, 179 classifiers on 121 datasets
\citep{fernandezdelgado2014hundreds}, 45 tabular benchmarks against deep
architectures \citep{grinsztajn2022tabular, shwartzziv2022tabular} and 41
classifiers on business credit data \citep{lessmann2015benchmarking} ---
boosted and bagged tree ensembles are the consistent top tier, boosting usually
ahead by a small margin. In e-commerce and logistics the same family is the
working choice: quantile regression forests for delivery times at JD.com
\citep{salari2022delivery}, tree models for late supplier deliveries
\citep{baryannis2019risk}, and random forests or boosted trees in the published
studies on this very dataset \citep{fouad2023satisfaction,
wangkiat2023satisfaction, gopinadh2025satisfaction}, none of which controls
leakage by prediction moment --- the gap Section~\ref{sec:data} closes.

The boosted learner used here is \code{HistGradientBoosting}: LightGBM's
histogram-binned design \citep{ke2017lightgbm} implemented inside
\code{scikit-learn} \citep{sklearn2024hgb}, so it stays within the brief's
toolset rather than adding XGBoost \citep{chen2016xgboost}; it handles missing
values natively, which the cold-start track-record features need; and it fits
an order of magnitude faster than a comparable forest. Heterogeneous ensembles
follow \citet{wolpert1992stacking}. The literature supplies the ranking; the
ladder of Section~\ref{sec:method} supplies what it cannot --- whether, on this
data, each rung beats the one below it.

\subsection{Feature engineering on a relational order model}

Deriving features from nine tables raises a problem a flat file does not: three
of them are one-to-many against orders, so a naive join multiplies rows and
reweights every aggregate towards multi-item orders, and each child relation
must be aggregated to the analysis grain before joining. The feature groups then
follow the literature. Origin--destination structure and distance are the
backbone of published delivery-time models, from JD.com's arrival forecasts
\citep{salari2022delivery} to last-mile travel-time predictors
\citep{liu2021lastmile}; the quoted promise encodes what the platform's routing
knew \citep{cui2024promise}; order size, value and freight are the order-level
controls through which \citet{deshpande2023logistics} identify the
logistics--rating link, and demand peaks stress fulfilment. The past performance
of the parties to a transaction is the standard predictor in supply-chain delay
models \citep{baryannis2019risk}, and an as-of aggregate of a key's earlier
outcomes is the time-respecting form of target encoding
\citep{miccibarreca2001target}, which must be fitted inside the
cross-validation loop. The general vocabulary of counts, ratios and calendar
encodings follows \citet{kuhn2019feature} and \citet{heaton2016feature}.

Leakage is the dominant risk. \citet{kaufman2012leakage} formalise it as
information about the target unavailable at prediction time and identify the two
forms that threaten this dataset, post-outcome features and improper
partitioning; \citet{kapoor2023leakage} show how routinely both survive peer
review, across 294 affected papers in 17 fields. Ours are their
\emph{illegitimate feature} and \emph{no clean separation} categories, and the
framing makes leakage relative to a \emph{prediction moment}: delivery outcomes
are forbidden when forecasting delivery and admissible when predicting a review
not yet written. On partitioning, \citet{bergmeir2012cv} show that standard
$k$-fold cross-validation is invalid for temporally ordered data.

\subsection{Evaluation under imbalance, drift and capacity}

Both classification targets are heavily imbalanced. \citet{saito2015prplot}
show that precision--recall curves are more informative than ROC curves under
strong imbalance, because ROC-AUC is insensitive to a large true-negative mass;
PR-AUC is reported here against its no-skill floor, with raw accuracy as a
diagnostic beside the majority-class rule. The obvious correction is one the
project tests rather than applies: \citet{goorbergh2022imbalance} find that
under-sampling, over-sampling and SMOTE \citep{chawla2002smote} do not improve
discrimination and badly miscalibrate the resulting probabilities, and recommend
adjusting the decision threshold on uncorrected data instead
(Section~\ref{sec:results-review} reports the test). Calibrated probabilities
are then a requirement, because a cost-sensitive decision multiplies them by a
cost \citep{elkan2001cost}, and calibration is reported, not assumed
\citep{niculescu2005calibration, platt1999probabilistic}.

Regression errors are judged against a na\"ive benchmark
\citep{hyndman2006accuracy}, decisive here because delivery performance improves
dramatically across the window: under such a level shift $R^2$ against a
baseline fitted on the past turns negative even for a useful model, whereas mean
absolute error stays interpretable; this is concept drift \citep{gama2014drift},
whose recommendation that non-stationary models be refitted with respect to
time shapes the deployment plan. A ranking metric summarises the whole list,
but a team contacts only the $K$ riskiest orders it can reach in a day, so the
quantity that describes its day is the precision and recall of that top-$K$ ---
the targeting view of \citet{provost2013data}, and the top-decile lift that
profit-driven evaluation uses in place of accuracy \citep{verbeke2012churn,
lessmann2015benchmarking}. Whether an intervention changes the outcome is a
causal question observational data cannot answer; the standard instrument is a
randomised controlled experiment \citep{kohavi2009controlled}, which
Section~\ref{sec:pilot} designs. The deployment follows the practice
literature: a model card stating what is served and where it fails
\citep{mitchell2019modelcards}, and an explicit account of the version and data
dependencies that quietly break production systems \citep{sculley2015debt}.
All modelling uses \code{scikit-learn} pipelines \citep{pedregosa2011sklearn},
keeping preprocessing inside each fold.

\subsection{The gap addressed}

The service-quality literature converges on fulfilment as the dominant
observable driver of e-commerce satisfaction but measures it one way while the
marketplace measures another; the tabular machine-learning literature converges
on tree ensembles, and the e-commerce applications that use them rarely say
\emph{when} the model would run; and the methodological literature is emphatic
that on imbalanced, temporally ordered data the choice of split, metric and
imbalance handling decides whether a number means anything. The contribution
here is narrow and concrete: make the prediction moment a first-class design
decision, measure its cost, and evaluate at the capacity the stakeholder has.
